\message{ !name(DiffCalc.tex)}
\message{ !name(InverseTrigDiff.tex) !offset(-2) }
\section{The Inverse Tangent and Cotangent Functions }
\label{sec:inverse-tan-and-cot}
\markright{\sc{}Inverse Tangent and Cotangent Functions }
% \baselineskip=1.5cm
% \vskip-.35in

If $x$ is the tangent of $y$, $x=\tan(y)$, then we say that $y$ is the
\term{arctangent} of $x$, $y=\arctan{x}$. Speaking loosely, we'd like
for the arctangent to be a  function which ``undoes'' the
tangent.

% But you have surely noticed that the information given in the
% table above presents some difficulties to
% computation. The presence of the ``$\pm$'' symbol suggests that when
% we compute the derivative of the arcsine, arccosine, arcsecant, and
% arccosecant functions we get to choose whether to choose positive or
% negative value.

% But this seems nonsensical, doesn't it? After all the slopes of each
% of these curves at a given value of $x$ will be positive or negative
% independent of any choices we might make. It is equally odd that
% this plus-or-minus conundrum doesn't seem to be a problem with the
% arctangent or the arccotangent functions. Why would that be?

% The problem here -- or at least part of the problem -- is that we
% haven't actually \emph{clearly defined} what we mean by
% ``$\arctan$,'' or any of the other ``arc'' functions. We will
% address this next.

\centerline{\includegraphics*[height=2.2in,width=5.5in]{../Figures/arctan3}}
%\centerline{\includegraphics*[height=2in,width=5in]{../Figures/arctan3}}
%\centerline{\includegraphics*[height=1.5in,width=4in]{../Figures/arctan3}}
The sketch above shows the graph of $y=\tan(x)$ (in red) and the same
graph (in blue) reflected about the graph of the line $y=x$ (dashed,
in black). Reflecting in this manner swaps the coordinates of each
point. For example, the point $(\pi, 0)$ is on the graph of
$y=\tan(x)$, because $\tan\left(\pi\right)=0$. Similarly, the point
$\left(0,\pi\right)$ is on the graph of $y=\arctan(x)$.

But is it true that $\arctan(0)=\pi$? If we are to believe our graph,
it could be true. But couldn't it also be true that
$\arctan(0)=2\pi$, or even $\arctan(0)=-\pi$?  The blue graph above is
clearly the graph of the arctangent as we have defined it but we seem
to have a choice for the output of $\arctan(x)$. Many choices, in fact.


Having more than one output for a given input clearly violates what we
mean when we use the word ``function.''  A mathematical function
returns \emph{exactly one} output for a given input. There is no
choice. The sketch of the arctangent (in blue) is not the graph of a
\term{function}. It is the graph of the \term{multifunction}
$\arctan(x)$. Each of the blue curves on the right is one
\term{branch} of the multifunction\aside{Note that by itself, any
  single branch of the arctangent \emph{is} a function.},
$\arctan(x)$. In fact all of the ``arc'' functions from trigonometry
are multifunctions.  Mult-functions are interesting objects and are
well worth studying.  But this is not the time for that study. Right
now we are only interested in the properties of the \term{inverse
  function} of $\tan(x)$. Since $\arctan(x)$ is not a function, it
can't be the inverse function of $\tan(x)$ so we will stop thinking
about it as soon as we can.

But if $\arctan(x)$ is not the inverse of $\tan(x)$ what is? Since we
only allow one output for a given input, could it
be as simple as choosing just one of the branches of the
$\arctan(x)$, say the one that lies between $y=-\frac{\pi}{2}$ and
$y=\frac{\pi}{2}$ and calling it the the inverse of  $\tan(x)$?\\
\centerline{\includegraphics*[height=2.7in,width=4.5in]{../Figures/arctan5}}
Actually, yes. It is just that simple. And the usual choice is the one
we've indicated.  In this text we will designate this branch of
$\arctan(x)$ with $\inverse\tan(x)$. Thus the symbol $\inverse\tan(x)$
denotes the \term{inverse tangent} function. Because the multifunction
arctangent is not a function we will henceforth ignore
it as much as possible.



\begin{Digression}[Inverse Function Notation]
\label{digression:inverse-function-notation}
      %       It is easy to understand what the inverse of a function is: If the
      %       function $\tan$ takes $x$ to $y$, then the function $\inverse\tan$
      %       ``undoes'' the function $\tan$ (If $y=\tan(x)$ then
      %       $x=\inverse\tan(y)$).

      %       The \emph{concept} really is that simple.


      %       But\notefrombob{I have also seen the notation y=Arctan x to denote the
      %       principal branch instead of y=x .  Do we want to utilize this at all
      %       in the discussion about inverse functions or is this complicating
      %       things unnecessarily? -- Jul 19, 2021} function inversion can be
      %       \emph{very} difficult to work with. This is partly because it is so
      %       abstract, but also because the notation we use can be very
      %       confusing. Probably the clearest notation we have is the ``arc''
      %       designation. The symbol ``$\arctan,$'' for example, clearly signals
      %       that the function\aside{In the computer age this has been shortened to
      %       ${\rm atan}$ for reasons that are no longer apparent.}, $\arctan(x)$
      %       returns an angle (the ``arc'' of a circle).


      %       We (the authors) have consistently used the ``arc'' prefix to
      %       designate the multifunctions which gave us the inverse trigonometric
      %       functions. Thus by our convention $\arctan(x)$ and $\inverse\tan(x)$ refer
      %       to different concepts altogether. We will continue this practice throughout this
      %       text. That is, we  will always, and only, refer to inverse functions by placing
      %       ``$-1$'' as an exponent on the name of the original function. Thus if
      %       $f(x)$ is an invertible function then $\inverse f(x)$ is it's inverse.

      %       We do this partly to avoid confusion in what can be a very confusing
      %       topic. But this is by no means a standard usage. 
      %       In fact for most scientists, engineers, and mathematicians
      %       ``$\arctan(x)$'' and $\inverse\tan(x)$, are synonymous.  In the
      %       computer age it has also become common to use $\text{atan}(x)$, and
      %       there are other variants as well.


The use of $-1$ as an exponent is probably the most common notation
used to indicate function inversion. Unfortunately, from the
standpoint of a student it is also probably the worst possible
notation we could have invented because It is so very easily confused
with reciprocation. While it is true that
  \begin{equation}
    \label{eq:NegOneRecip}
    2^{-1} = \frac12,
\end{equation}
it is important to remember that
\begin{equation}
  \label{eq:NegOneInv}
  \inverse{\tan}(x) \text{ is \emph{not} equal to } \frac{1}{\tan(x)}.
\end{equation}

%   \begin{wrapfigure}[21]{r}{3.2in}
%     \vskip-6mm{}
%     \captionsetup{labelformat=empty}
% %    \centerline{\href{https://mathwithbaddrawings.com/}{\includegraphics*[height=1.5in,width=2.2in]{../Figures/BadDrawing-InvSine}}}
% %    \centerline{\href{https://mathwithbaddrawings.com/}{\includegraphics*[height=2.5in,width=3.2in]{../Figures/BadDrawing-InvSine}}}
%     \centerline{\href{https://mathwithbaddrawings.com/}{\includegraphics*[height=2in,width=3.2in]{../Figures/BadDrawing-InvSine}}}
%     \label{fig:BadDrawing-InvSine}
%     \caption{\textcolor{black}{\centerline{The World's Most Evil Mathematician}
%         \centerline{from {Math
%             With Bad Drawings}}
%           \centerline{by Ben Orlin}}}
%   \end{wrapfigure}
  That the ``$-1$'' notation is
  used for both comes from the fact that both $2^{-1}$ and
  $\inverse\sin(x)$ really are inverses. But they are different kinds
  of inverses: $\inverse2$ is the \emph{multiplicative} inverse of
  $2$, which means simply that
  $$\frac12\cdot2=1 \text{ and that } 2\cdot\frac12 =1.$$ 
  On the other hand $\inverse\tan(x)$ is the \emph{functional} inverse of
  $x=\tan(y)$, which means that if
  $-\frac{\pi}{2}\lt{}x\lt{}\frac{\pi}{2}$ then
  $$
  \inverse\tan(\tan(x)) = x, 
  $$
  and that
  $$
  \tan(\inverse\tan(x))=x.
  $$

  In equation~\ref{eq:NegOneRecip} you can think of the $-1$ in the
  exponent as an operator. It operates on $2$ by taking its
  reciprocal: $\inverse{2}=\frac12$. In equation~\ref{eq:NegOneInv}
  the $-1$ is not an operator. It is part of the symbol,
  $\inverse\tan$, that we use to denote the inverse tangent.
  
  With practice this all gets easier, but in the beginning it is very
  troublesome. Be careful.

  \begin{ProficiencyDrill}
    Evaluate each of the following:
           \begin{enumerate}[label={\bf{}(\alph*)}]{}
        \begin{multicols}{4}
         \item $\inverse{\left(\frac{2}{\inverse{3}}\right)}$
         \item $\inverse\tan\left(\frac{\sqrt{3}}{2}\right)$
         \item $\inverse{\left(\cot\left(\frac{\pi}{3}\right)\right)}$
         \item $\inverse{\left(\inverse\tan\left(1\right)\right)}$
        \end{multicols}
       \end{enumerate}
%    \vskip-4mm{}
     \end{ProficiencyDrill}

\end{Digression}

But what, exactly, is $\inverse\tan(x)$ the inverse of? From the sketch above it
should be clear that restricting the \alert{range} of
$\inverse\tan(x)$ to values in the interval $(-\pi/2, \pi/2)$ forced
us to restrict the \alert{domain} of $\tan(x)$. So
$$
\inverse\tan(x) \text{ with domain all real numbers,} \RR
$$
is the inverse of
$$
\tan(x) \text{ with domain, } (-\pi/2, \pi/2),
$$
and it is \emph{not} the inverse of
$$
\tan(x) \text{ with domain, } \RR
$$
because that function  has no inverse.


An easy way to remember what $\inverse{\tan}(x)$ means is to read the
symbol $\inverse{\tan}(x)$ as\\
\centerline{``the angle whose tangent is $x$.''}
The
advantage of this phrasing is that it emphasizes that $x$ is the
tangent of some angle, and $\inverse{\tan}(x)$ is that
angle. Similarly the inverse functions of the other trigonometric
functions should be
read as\\
\centerline{``the angle whose \{sine, cosine, secant, whatever\} is
  $x$''.}  Each of these will also come with suitable restrictions on
its range.

We should point out that reserving the ``arc'' notation for the
trigonometric multifunctions is actually a conceit of the authors of
this text. Most of the world uses $\arctan(x)$ and $\inverse\tan(x)$
interchangeably, and eventually you will too. But for now it will be
helpful for you to keep in mind that the difference between them is
that restricting the range  of   $\inverse\tan(x)$  to
$-\frac{\pi}{2}\lt x\lt\frac{\pi}{2}$ guarantees that it is a
function. But since the range of   $\arctan(x)$ is not similarly
restricted it is not
a function.


\begin{ProficiencyDrill}
  Suppose we had chosen
  $\frac{\pi}{2}\lt{}\inverse\tan(x)\lt{}\frac{3\pi}{2}$. What function (with
  domain) is that the inverse of?
\end{ProficiencyDrill}

      %       The figure above shows the graph of $\tan(x)$ and its inverse,
      %       $\inverse\tan(x)$. But notice that we've explicitly constrained
      %       the domain of
      %       $\tan(x)$. Do you see why we had to do that?  We can only include one
      %       branch of $\tan(x)$ because $\inverse\tan(x)$ has only one branch of
      %       $\arctan(x)$ and they have to match up if they are going to be
      %       mutually inverse.


      %       Each of the curves in this graph is called a branch of the arctangent.
      %       The graph in red above is called the \term{principle
      %       branch} because, by convention, it is the one chosen by default --
      %       unless there is a compelling reason to choose another.  This is the one where
      %       $-\frac{\pi}{2}\lt{}y\lt{}\frac{\pi}{2}$.one.

All of this fussiness is really just about making our abstract
definitions useful  and
consistent. It would be nice if these details only impinged on us in an
abstract setting but unfortunately some  practical difficulties do
come up, as the following drill shows.

%\pagebreak[4]
% \begin{ProficiencyDrill}
%   Compare the graphs  of each of the following:
%   \begin{enumerate}[label={\bf{}(\alph*)}]
%     \begin{multicols}{4}
%     \item $y=\inverse\tan\left(\frac1x\right)$
%     \item $y=\inverse\tan\left(\frac{1}{x-1}\right)$
%     \item $y=\inverse\tan\left(\frac{1}{x-1}\right)$
%     \item $y=\inverse\tan\left(\frac{x+1}{1-x}\right)$
%     \end{multicols}
%   \end{enumerate}
%   % $y=\inverse\tan\left(\frac{1+x}{1-x}\right)$ to the graph of
%   % $y=\inverse\tan(x)+\inverse\tan(1)=\inverse\tan(x)+\frac{\pi}{4}$.\\
%   % \centerline{\includegraphics*[height=1in,width=4in]{../Figures/arctanabs1}}
%   Can you explain why they are discontinuous even though $y=\inverse\tan(x)$
%   is continuous?
% \end{ProficiencyDrill}

\begin{embeddedproblem}
  \vskip-2mm{}
  Find \emph{all} solutions of  $\tan(x)=1$ and
  $x=\inverse\tan\left(1\right)$. Do they have the same set of
  solutions?\\
  \hint{Obviously, they do not. Otherwise we wouldn't have asked the
    question. What is the difference between the two sets of
    solutions?}
\end{embeddedproblem}

\begin{Digression}[\boldmath{}$\tan(x)$ Has No Inverse]
  The function $\tan(x)$ is not invertible. This is because, by
  custom, its domain is all real numbers, $\RR$. Thus it is
  \alert{not} the same function as
  $$\tan(x) \text{ with domain } -\frac{\pi}{2}\lt{}x\lt{}\frac{\pi}{2}.$$  A
  function has two parts: (1) The set of inputs (domain) and (2) the
  rule associating input with output. Most of the time the domain is
  the set of all real numbers ($\RR$), so we don't bother to
  explicitly state that the domain of, say $f(x)=x^2$ is
  $\RR${}. Unless otherwise specified we simply assume that it is.
  Properly speaking, we should \emph{always} specify the domain of our
  functions, but typically we don't. The custom is to assume that the
  domain of a function is $\RR$ unless otherwise stated.
  % That is confusing enough, but we (mathematicians) add to the confusion when dealing with
  % inverse functions because we habitually violate this custom as well.
  Thus the  function, $\tan(x)$ cannot be inverted because on its
  domain ($\RR$) it has more than
  one branch.

  Because we both want and need an inverse tangent function we
  restrict the range of $\inverse\tan(x)$. But this is the inverse of
  $\tan(x)$, with domain $-\frac{\pi}{2}\lt{}x\lt{}\frac{\pi}{2}$,
  which has only one branch. It is \alert{not} the inverse of
  $\tan(x)$, because it has too many branches.

  To use our notation as it was intended to be used we \emph{should}
  say that $\inverse\tan(x)$ is the inverse of $\tan(x)$, with domain
  $-\frac{\pi}{2}\lt{}x\lt{}\frac{\pi}{2}$, which is true. But instead, we
  typically just say that it is the inverse of $\tan(x)$, which
  strictly speaking, is not true. 


  When we (mathematicians) talk amongst ourselves this is not a
  problem. We all understand what we mean. But it can be very
  confusing for students who usually have a more tenuous grasp of
  these distinctions.  It is a bit unfair of us (mathematicians) to
  speak to students in what amount to incomplete sentences, but we
  tend to do it anyway, mostly out of habit.

  We (the authors) apologize.
\end{Digression}

The definition of the arccotangent is similar to the definition of the
arctangent: If $x$ is the cotangent of $y$, $x=\cot(y)$, then we say
that $y$ is the \term{arccotangent} of $x$, $y=\arccot(x)$. The
arccotangent is also a multifunction with multiple branches so, just
as before, we will have to decide which branch to use to define the
inverse cotangent: $\inverse\cot{x}$.


We obtained our branch for the single-valued function $y=\inverse\tan(x)$
from the multifunction $y=\arctan(x)$ by restricting its  range to
$-\frac{\pi}{2}\lt y\lt\frac{\pi}{2}$. A similar restriction would
allow us to obtain the single-valued
function $\inverse\cot(x)$ from the multifunction $\arccot(x)$, but it
is a bit simpler to once again use a trigonometric  identity.

\begin{ProficiencyDrill}
  Use the identity $$\inverse\cot(x) = \frac{\pi}{2}-\inverse\tan(x)$$
  to choose an appropriate (define) domain of $y=\inverse\cot(x)$.
\end{ProficiencyDrill}


\section{The Witch of Agnesi and Inverse Tangent Function }
\label{sec:witch-agnesi-inverse}
\begin{wrapfigure}[]{l}{1.7in}
  \captionsetup{labelformat=empty}
  \centerline{\href{https://mathshistory.st-andrews.ac.uk/Biographies/Agnesi/}{\includegraphics*[height=2in,width=1.7in]{../Figures/AgnesiPortrait}}}
  \caption{\centerline{Maria Ga\"{e}tana Agnesi}\\ \centerline{(1718-1799)}}
  \label{fig:AgnesiPortrait}
\end{wrapfigure}
In a time when women had very few options in life
\href{https://mathshistory.st-andrews.ac.uk/Biographies/Agnesi/}{Maria
  Ga\"{e}tana Agnesi} (1718-1799) was both exceptional and very
lucky. Her brilliance and talent were recognized early and nurtured by
her wealthy father, who encouraged her studies and provided her with
the best possible tutors to develop her talents. She mastered the
Calculus of Newton and Leibniz and in $1748$ wrote a series of
textbooks on the topic titled \emph{\href{https://www.maa.org/press/periodicals/convergence/mathematical-treasures-maria-agnesis-analytical-institutions}{Instituzioni Analitiche ad Uso
  Della Giovent\`u{} Italiana}\aside{This is usually translated as
    \pubtitle{Analytical Methods for the use of Italian Youth}}.} It
was immediately recognized as a masterpiece of mathematical
exposition and  was used throughout Europe during eighteenth century. It
was by far the most popular Calculus textbook in use at the time.  An
Englishman,
\href{https://mathshistory.st-andrews.ac.uk/Biographies/Colson/}{John
  Colson}, was so impressed with Agnesi's work that in $1760$ he took it
upon himself to learn Italian specifically so that he could translate
her text into English.

In her text Agnesi\aside{Her name is pronounced ``Awn-yes-ee,'' not
  ``Ag-nes-ee.'' } collected many of the known results of the time and
organized them for the instruction of students. The Latin name of one
of the curves that she used for instruction is \emph{versoria} (``rope
that turns a sail'') because of its shape. Agnesi correctly translated
this into Italian as ``la versiera.''  Unfortunately Colson mistook
this for ``l'aversiera'' which means ``witch.'' As a result this
particular curve has been known ever since as the Witch of Agnesi.


\begin{wrapfigure}[13]{r}{4.4in}
  \vskip-6mm{}
  \captionsetup{labelformat=empty}
  \includegraphics*[height=2.2in,width=4.4in]{../Figures/Witch1}
\label{fig:Witch1}
\end{wrapfigure}
As modern function notation was yet to be invented Agnesi
understood this curve geometrically as a particular set of
points. In the  diagram at the right think of $P$ as moving from left
to right. For each value of $x_0$, draw the line
from the origin to $P$ and locate the vertical
coordinate $y_0$ where this line intersects the circle
$$x^2+ \left( y-\frac12 \right)^2=\left(\frac12\right)^2,$$
 centered at
 $\left(0, \frac12\right)$, with radius $r=\frac12$.
 Every
point $(x_0,y_0)$ is a point on the Witch of Agnesi. The Witch itself  is the
curve shown in red.

%\pagebreak[4]
\begin{embeddedproblem-1line}{}
  In the diagram above show that the coordinates $x_0$ and $y_0$
  satisfy the equation, $y_0=\dfrac{1}{1+x_0^2}$. Thus the Witch is
  the graph of
  $$
  y(x)=\frac{1}{1+x^2}.
  $$
% \hint{Notice that $(x_0, y_0)$  is not the point of intersection of the circle and the line, but they do have the
% same $y$ coordinate. It might be  helpful to label this intersection
% point $(x_1, y_0)$ and notice  that you are
% trying to find $y_0$.}
%  \vskip-6mm{}
\end{embeddedproblem-1line}

%%% \centerline{  \includegraphics*[height=2.5in,width=4in]{../Figures/Witch2}}

Whenever you encounter a new curve an important question to ask is,
``What is its derivative?'' In this case you already have everything you
need to find the answer.

\begin{embeddedproblem-1line}
  Show that if $y=\frac{1}{1+x^2},$ then
  $\dfdx{y}{x}=\frac{-2x}{(1+x^2)^2}.$
\end{embeddedproblem-1line}

\begin{wrapfigure}[10]{r}{3in}
  \vskip-2mm{}
\captionsetup{labelformat=empty}
\centerline{\includegraphics*[height=1.5in,width=3in]{../Figures/Witch3}}
\label{fig:Witch3}
\end{wrapfigure}
Computing the derivative of a function as a formula is very
useful. This has in fact been our primary focus so far in this
text. However to truly understand the relationship between a function
and its derivative nothing can replace seeing both of them graphed
together.  Look closely at the relationship between the Witch and its
derivative in the sketch at the right. Is it clear how these graphs
are related?

Notice in particular that when $x=0$ the $y$ coordinate of the Witch
is at its highest point (maximum value), whereas the $y$ coordinate of
$\dfdx{(\text{Witch})}{x}$ is zero.



% \begin{wraptable}{r}{2.4in}
%   \captionsetup{labelformat=empty}
%   \begin{tabular}{ |c|c| }
      %       \hline
      %       \multicolumn{2}{|c|}{} \\
      %       \multicolumn{2}{|c|}{\Large\bf{}Differentiation of } \\
      %       \multicolumn{2}{|c|}{} \\
      %       \multicolumn{2}{|c|}{\Large\bf{} Inverse Trig Functions} \\
      %       \multicolumn{2}{|c|}{} \\
      %       \hline
      %       $y$   \amp{}$\dx{y}$\\\hline\hline
      %       $\inverse\sin (x)$ \amp{}    $\dfrac{1}{\sqrt{1-x^2}}\dx{x} $\bigstrut{}\\
      %       $\inverse\cos (x)$ \amp{} $\dfrac{-1}{\sqrt{1-x^2}}\dx{x}$ \bigstrut{}\\
      %       $\inverse\tan (x)$\amp{}$\dfrac{1}{1+ x^2}\dx{x}$\bigstrut{}\\
      %       $\inverse\cot (x)$\amp{}$\dfrac{-1}{1+x^2}\dx{x}$\bigstrut{}\\
      %       $\inverse\sec (x)$\amp{}$\dfrac{1}{\abs{x}\sqrt{x^2-1}}$\bigstrut{}  \\
      %       $\inverse\csc x$\amp{}$\dfrac{-1}{\abs{x}\sqrt{x^2-1}}\dx{x}$\bigstrut\\\hline
      %     \end{tabular}
      %       \caption{\textcolor{blue}{\bf{} Memorize these!}}
      %       \label{tab:InvTrigDiffRules}
      %       \end{wraptable}
      % %     \end{center}
We know that the derivative of a curve at a given point gives us the
slope of the line tangent at that point. Thus at its highest point the
line tangent to the Witch will be horizontal.  That is, its derivative
will be zero. This is exactly what these two graphs are showing us.

When $x$ is just to the right of zero, the slope of the Witch is close
to zero and negative. Similarly just to the left of zero the slope of
the Witch is close to zero and positive. At the extreme right end of
the graph above the slope of the Witch is close to zero and negative
and at the extreme left it is close to zero and positive. And all of
this is reflected in the shape of the blue (derivative) curve.

It is a little more interesting to reverse this process. That is,
suppose we could see only the derivative of the Witch (blue
curve). Could we figure out the shape of the Witch from this?

Sure we can. Well, almost.

Consider: From the extreme left of the graph to $x=0$ the
$y-$coordinate of the (blue) derivative curve is above the $x-$axis. Therefore
the slope of the Witch is positive. Clearly wherever the slope of a
function is positive the function is increasing, so we can deduce that
the Witch increases from left to right until we reach $x=0$ where the
blue derivative function crosses the $x$-axis.

Thereafter the $y-$coordinate of the blue derivative curve is below the
$x-$axis. That is, the slope of the Witch is negative, so the Witch is
decreasing.

Just by looking at the derivative curve we can see that the Witch of
Agnesi increases from left to right until we reach $x=0,$ and after
that it decreases. Although this shows us the shape of the Witch it is
not enough to completely describe the Witch of Agnesi. We can tell a
lot about a curve by looking at the graph of its derivative but we
can't tell everything.

\begin{ProficiencyDrill}{}
  \label{drill:NonUniqueAntiderivative}
  In the sketch below the red curve is the graph of the Witch and the
  blue curve is the graph of its derivative. Convince yourself that
  the blue curve could also be the graph of the derivative of any of
  the other curves shown as well. As clearly as you can, explain what this
  suggests about the relationship between a
  function and its derivative.\\
  \centerline{\includegraphics*[height=1.5in,width=3in]{../Figures/Witch11}}
  
\end{ProficiencyDrill}


Next we'd like to apply this same sort of reasoning to the Witch
itself.  We'd like to answer the question, ``What curve is the Witch
of Agnesi the derivative of?''  Whatever that curve is, we'll call it
the \term{antiderivative} of the Witch for obvious reasons.

%\centerline{\includegraphics*[height=1.5in,width=3in]{../Figures/Witch6}}
\centerline{\includegraphics*[height=1.5in,width=3in]{../Figures/Witch6}}
As before, by simply looking at the graph of the Witch we can see
that: (1) At the extreme ends of its graph the antiderivative will
have a \emph{positive} slope which gets closer and closer to zero as
we go farther from the origin, but (2) since the Witch never crosses
the $x$-axis the antiderivative will never have a horizontal slope,
and (3) at $x=0$ the slope of the antiderivative will be equal to one. A sketch
of the antiderivative of the Witch begins to emerge when we put
these three observations into our graph:\\
\centerline{\includegraphics*[height=1.5in,width=3in]{../Figures/Witch7}}

Continuing to fill between the black segments in the same  fashion
the sketch becomes even clearer:\\ 
\centerline{\includegraphics*[height=1.8in,width=3in]{../Figures/Witch8}}
Finally, connecting the dots with a smooth curve we see the the
graph of the antiderivative of the Witch of Agnesi must look like this:\\
\centerline{\includegraphics*[height=1.75in,width=3.6in]{../Figures/Witch9}}

We've been a little lazy with our use of language.  As we saw in
Drill~\ref{drill:NonUniqueAntiderivative}, the antiderivative is not
unique, so it is improper to speak of \emph{the} antiderivative. Any
one of the curves shown below could be \emph{an} antiderivative of the
Witch of Agnesi. (Except the black one, of course. That's the Witch.)\\
\centerline{\includegraphics*[height=1.8in,width=3in]{../Figures/arctan1}}

This graph should look very  familiar to you. It looks very much
like the  graph of the arctangent function we discussed
in Section~\ref{sec:inverse-tan-and-cot}, doesn't it? Do you suppose
it is possible that the derivative of the
arctangent function is the Witch of Agnesi?

Of course it is. Why
  else would we have led you down this path? But there are some
  subtleties here that we shouldn't ignore. For example we will want
  to find a function that is the antiderivative of the Witch. We will proceed
  cautiously.

  \begin{embeddedproblem}
    Sketch the graph of the derivative and the antiderivative of each
    of the following curves on the same set if axes,
           \begin{enumerate}[label={\bf{}(\alph*)}]{}
        \begin{multicols}{2}
         \item\ \\ \leftline{\includegraphics*[height=1.5in,width=1.5in]{../Figures/DiffSketch1}}
         \item\ \\ \leftline{\includegraphics*[height=1in,width=1.5in]{../Figures/DiffSketch2}}
         \item\ \\ \leftline{\includegraphics*[height=1.5in,width=1.5in]{../Figures/DiffSketch3}}
         \item\ \\ \leftline{\includegraphics*[height=.8in,width=1.5in]{../Figures/DiffSketch4}}
        \end{multicols}
      \end{enumerate}
%      \vskip-10mm{}
  \end{embeddedproblem}

\section{The Differentials of the Inverse Tangent and Inverse
  Cotangent Functions}
\label{sec:diff-inverse-tancot}
\markright{Tangent and Cotangent Differentials}



\begin{wrapfigure}[9]{r}{3.5in}
  \vskip-3mm{}
\captionsetup{labelformat=empty}
\centerline{\includegraphics*[height=1.3in,width=3.5in]{../Figures/arctan6}}
\label{fig:arctan6}
\end{wrapfigure}
% When you learned
% Trigonometry you would have encountered the inverse  functions
% % ,
% % $\inverse\sin(x)$, $\inverse\cos(x)$, $\inverse\tan(x)$,
% % $\inverse\cot(x)$, \\$\inverse\sec(x)$, $\inverse\csc(x)$,
% of all six of
% the trigonometric functions. In particular
% the graph of $y=\inverse\tan(x)$ at the right  looks so much
% like the derivative of the Witch of Agnesi that it is tempting to
% simply assume that this is true.
Since we want to find a function whose derivative is the Witch of
Agnesi and we are pretty sure that each of the branches of
the arctangent multifunction takes the shape of the Witch's derivative it
is tempting to assume that the branch of the arctangent shown at the
right is the function we are looking for.

But we must be careful.  Reasoning from pictures, as we've just done,
can be very useful.  but the most we can hope for when reasoning from
pictures is an intuitive feel for the problem. We always need to
confirm our intuition analytically. Always.

So we will verify, analytically, our conjecture that the derivative of
the inverse tangent function is the Witch of Agnesi.

From our definition of $\inverse\tan(x)$ we know that
$
y=\inverse\tan(x)$ precisely when\aside{Notice how cavalierly we've
  ignored the \emph{necessary} constraints on the domain of the
  tangent. The function $\inverse\tan(x)$ is not the inverse of
  $\tan(x)$ because the domain of $\tan(x)$ is $\RR$. It is the
inverse of $\tan(x)$, with domain, $\frac{-\pi}{2}\lt x
\lt\frac{\pi}{2}$. The domain is part of the definition of the
function.} $\tan(y)=x$.
Differentiating $\tan(y)=x$ gives:
\begin{align*}
  \dx(\tan(y)) \amp{}= \dx{x}\\
  \sec^2(y)\dx{y} \amp{}= \dx{x}\\
  \dx{y}\amp{}= \frac{1}{\sec^2(y)}\dx{x}.
\end{align*}
This is a correct formula for the differential of the arctangent
function. However, there are two problems: (1) it does not seem to
match our conjecture, and (2) even if it is correct this formula will
not be very useful to us in this form since we have $\dx{y}$
in terms of $y$ itself. We can address both of these problems by
finding $\dx{y}$ in terms of $x$ and $\dx{x}$.

\begin{wrapfigure}[]{r}{2.3in}
  \vskip-6mm{}
  \captionsetup{labelformat=empty}
  \centerline{\includegraphics*[height=2.3in,width=2.3in]{../Figures/arctanMnemonic2}}
  \label{fig:arctanMnemonic2}
\end{wrapfigure}
Referring to the diagram at the right, recall from
Section~\ref{subsec:trig-interl} that $\tan(y)=\frac{x}{1}=x$ so
$y=\inverse\tan(x)$. The hypotenuse in the diagram at the right is
$\sec(y)$ so from the Pythagorean Theorem we see that
$\sec^2(y)=1+x^2$.
Thus
$$
\dx{y}=\frac{1}{\sec^2(y)}\dx{x}=\frac{1}{1+x^2}\dx{x},
$$
or
$$
\dfdx{y}{x}=\frac{1}{1+x^2}
$$
So our intuition was correct. The Witch is indeed the derivative of
the arctangent function.  A similar derivation can be used for the
arccotangent.

%\pagebreak[4]

\begin{embeddedproblem-enumerate}{}
  \begin{enumerate}[label={\bf{}(\alph*)}]
  \item The inverse cotangent function is defined as
    $y=\inverse\cot(x)$ if and only if $\cot(y)=x$.  Proceed as we did
    above to show that
    $$
    \dx{\left(\inverse\cot(x)\right)} = \frac{-1}{1+x^2}\dx{x}.
    $$
  \item Derive the same result from  the identity
    $\inverse\cot(x)=\pi/2-\inverse\tan(x)$.
  \end{enumerate}
\end{embeddedproblem-enumerate}

%\pagebreak[4]
\begin{embeddedproblem-enumerate}
  \begin{enumerate}[label={\bf{}(\alph*)}]
  \item Show that if $y=\inverse\tan(x)$ then $y$ satisfies the
    differential equation
    $$
    (1+x^2)\dfdxn{y}{x}{2}+2x\dfdx{y}{x}=0.
    $$
  \item The function   $y=\inverse\cot(x)$ satisfies the same differential
    equation. Show this in two different ways.
    \begin{enumerate}[label={\bf{}(\roman*)}]{}
    \item By direct computation, just as you did part (a).
    \item By direct computation, after first observing that $\inverse\cot(x)=\frac\pi2-\inverse\tan(x)$.
    \end{enumerate}
  \end{enumerate}{}
\end{embeddedproblem-enumerate}


\begin{embeddedproblem}
  Show that each of the following statements is true.
  \begin{enumerate}[label={\bf{}(\alph*)}]
    \begin{multicols}{2}
    \item $\dx{\left(\inverse\tan\left(\dfrac{1+x}{1-x}\right)\right)} = \dfrac{1}{1+x^2}\dx{x}$
    \item $\dx{\left(\inverse\tan\left(\dfrac{2+x}{1-2x}\right)\right)} = \dfrac{1}{1+x^2}\dx{x}$
    \item $\dx{\left(\inverse\tan\left(\dfrac{10+x}{1-10x}\right)\right)} = \dfrac{1}{1+x^2}\dx{x}$
    \item $\dx{\left(\inverse\tan\left(\dfrac{2543+x}{1-2543x}\right)\right)}
      = \dfrac{1}{1+x^2}\dx{x}$
    \end{multicols}
  \item Now show that
    $$\dx{\left(\inverse\tan\left(\dfrac{\kappa+x}{1-\kappa x}\right)\right)}
    = \dfrac{1}{1+x^2}\dx{x}$$ where $\kappa$ is any constant.
  \item This is weird isn't it? None of the functions on the left is
    $\inverse\tan(x)$  but they all have the same derivative as $\inverse\tan(x)$. Can you
    explain this?\\
    \hint{Let $y_1=\inverse\tan(x)$ and $y_2=\inverse\tan(\kappa)$. What is
      $\tan(y_1+y_2)$?}
  \end{enumerate}
%  \vskip-6mm{}
\end{embeddedproblem}



% We now have the differentials of all of the inverse trigonometric
% functions as seen in the table below.  The need for the domain
% restrictions in the center column should be clear if you think about
% the nature of the trig\-on\-o\-met\-ric functions. For example, since
% the sine (and cosine) function only return numbers between $-1$ and
% $1$ it makes no sense to write, for example, $\arcsin(2)$ because this
% asks for ``the arc (angle) whose sine is $2$,'' and there is no such
% arc\aside{At least not among the real numbers. It is possible to find
%   a complex ``angle'' whose sine is $2$. But this is well beyond the
%   scope of this text.}.  Similarly the inverse secant and inverse
% cosecant functions are only defined for values of $x\le -1$ and
% $x\ge1$ because the secant and cosecant functions only return values
% in that range.

% % \begin{wraptable}[21]{r}{4.3in}
% %   \vskip-1mm{}
% %   \captionsetup{labelformat=empty}
% \begin{center}
%     \begin{tabular}{|ccc|} \hline Function\amp{}
%       Domain\amp{}Derivative\\\hline\hline $y=\arctan(x)$ \amp{} Any
%       real
%       number\amp{}$\displaystyle\dfdx{y}{x}=\frac{1}{\sec^2(y)}=\frac{1}{1+x^2}$\\
%       \amp{}\amp{}\\
%       $y=\arccot(x)$ \amp{} Any real number\amp{}$\displaystyle\dfdx{y}{x}=\frac{-1}{\csc^2(y)}=\frac{-1}{1+x^2}$\\
%       \amp{}\amp{}\\
%       $y=\arcsin(x)$ \amp{} $-1\le x\le 1$\amp{}$\displaystyle\dfdx{y}{x}=\frac{1}{\cos(y)}=\frac{\pm1}{\sqrt{1-x^2}}$\\
%       \amp{}\amp{}\\
%       $y=\arccos(x)$ \amp{} $-1\le x\le 1$\amp{}$\displaystyle\dfdx{y}{x}=\frac{-1}{\sin(y)}=\frac{\mp1}{\sqrt{1-x^2}}$\\
%       \amp{}\amp{}\\
%       \multirow{2}{*}{$y=\arcsec(x)$} \amp{} $x\le-1$\amp{}\multirow{2}{*}{$\displaystyle\dfdx{y}{x}=\frac{1}{\sec(y)\tan(y)}=\frac{\pm1}{x\sqrt{x^2-1}}$}\\
%       \amp{}or $x\ge1$\amp{}\\
%       \amp{}\amp{}\\
%       \multirow{2}{*}{$y=\arccsc(x)$} \amp{} $x\le-1$\amp{}\multirow{2}{*}{$\displaystyle\dfdx{y}{x}=\frac{-1}{\csc(y)\cot(y)}=\frac{\mp1}{x\sqrt{x^2-1}}$}\\
%       \amp{}or $x\ge1$\amp{}\\
%       \amp{}\amp{}\\\hline
%     \end{tabular}
%   \end{center}

% \caption{\textcolor{blue}{Differentials of Inverse Trigonometric functions}}
%    \label{tab:InvTrigDiff1}
% \end{wraptable}



\section{The Other Inverse Trigonometric Functions}
\label{sec:other-inverse-trig}
Despite how much time we spent talking about them, the single valued
branches for $y=\inverse\tan(x)$ and $y=\inverse\cot(x)$, were
actually fairly apparent from the graphs of $x=\tan(y)$ and
$x=\cot(y)$. All we had to do was choose one continuous branch from
either graph. We described the process carefully because for the
other inverse trigonometric functions things are a little more
problematic. For those we will have to use the derivatives of our
``arc'' multifunctions as a guide, rather than the multifunctions
themselves.

In general, we define an arc[function] as $\text{arcfunction}(x)=y$
precisely when $\text{function}(y)=x$. So, $\arcsin(x)=y$ when
$\sin(y)=x$ and $\arccos(x)=y$ when $\cos(x)=y$.

\subsection*{The Inverse Sine Function: \boldmath{}$\inverse\sin(x)$}
\begin{wrapfigure}[12]{r}{1.5in}
  \vskip-5mm{}
  \captionsetup{labelformat=empty}
  % \centerline{\includegraphics*[height=1.9in,width=1.3in]{../Figures/arcsin5}}
  \centerline{\includegraphics*[height=2.2in,width=1.5in]{../Figures/arcsin5}}
\label{fig:arcsin5}
\end{wrapfigure}
When we graph $\arcsin(x)$ by flipping the graph of $\sin(x)$ across
the graph of the line $y=x$ we get the sketch at the right. For the
value of $x$ shown in the sketch the arrows point to a few of the
possible values of $\arcsin(x)$. There are infinitely many, so
obviously $\arcsin(x)$ is also a multifunction, not a function.


To find the  differentials of $y=\arcsin(x)$ and $y=\arccos(x)$ we
proceed just as we did with the $\arctan(x)$ and $\arccot(x)$. If
$y=\arcsin(x)$ then $x=\sin(y)$ so:
\begin{align*}
  \dx{(\sin(y))}\amp{}= \dx{x}\\
  \cos(y)\amp{}=\dx{x}\\
  \dx{y}\amp{}=\frac{1}{\cos(y)}\dx{x}
\end{align*}

\begin{ProficiencyDrill}
  \label{drill:InvSin1}
  Use the  diagram  below,
  \\
  \centerline{\includegraphics*[height=1in,width=1.6in]{../Figures/CosInvSin1}}
   to show that
   \begin{equation}
     \label{eq:InvSin11}
     \dx(\arcsin(x))=\frac{\pm 1}{\sqrt{1-x^2}}.
\end{equation}
\end{ProficiencyDrill}

%   \begin{wrapfigure}[21]{r}{3.2in}
%     \vskip-6mm{}
%     \captionsetup{labelformat=empty}
% %    \centerline{\href{https://mathwithbaddrawings.com/}{\includegraphics*[height=1.5in,width=2.2in]{../Figures/BadDrawing-InvSine}}}
% %    \centerline{\href{https://mathwithbaddrawings.com/}{\includegraphics*[height=2.5in,width=3.2in]{../Figures/BadDrawing-InvSine}}}
    % \centerline{\href{https://mathwithbaddrawings.com/}{\includegraphics*[height=2in,width=3.2in]{../Figures/BadDrawing-InvSine}}}
    % \centerline{The World's Most Evil Mathematician}
    %     \centerline{from {Math
    %         With Bad Drawings}}
    %       \centerline{by Ben Orlin}
  %   \label{fig:BadDrawing-InvSine}
  %   \caption{\textcolor{black}{\centerline{The World's Most Evil Mathematician}
  %       \centerline{from {Math
  %           With Bad Drawings}}
  %         \centerline{by Ben Orlin}}}
  % \end{wrapfigure}

\begin{embeddedproblem-enumerate}{}
  \label{drill:InvCos1}
  \begin{enumerate}[label={\bf{}(\alph*)}]
  \item Use the fact that $y=\arccos(x)$ if and only if $\cos(y)=x$
    and proceed as we did in Drill~\#\ref{drill:InvSin1} to show that
    \begin{equation}
      \label{eq:InvCos1}
      \dx\left(\arccos(x)\right)= \frac{\mp 1}{\sqrt{1-x^2}}\dx{x}.
  \end{equation}
    \comment{Notice how this formula differs from the one you found
      in Drill~\ref{drill:InvSin1}.}
  \item Derive the same result from the identity:  $\arccos(x)=
    \frac{\pi}{2}-\arcsin(x)$.
  \end{enumerate}
\end{embeddedproblem-enumerate}

The plus/minus and minus/plus that appear in
equations~\ref{eq:InvSin11} and~\ref{eq:InvCos1}, respectively,
are present because, $\arcsin(x)$ and $\arccos(x)$ are both
multifunctions like $\arctan(x)$.

Since $\arcsin(x)$ is not a function at all it can't be the inverse
\emph{function} of $\sin(x)$.  Once again, to \emph{create} an inverse
sine function (which we will call $\inverse\sin(x)$) we will have to
restrict the range of $\inverse\sin(x)$ just like we restricted the
range of $\inverse\tan(x)$.  The question is, what restriction should
we impose?

    \centerline{\href{https://mathwithbaddrawings.com/}{\includegraphics*[height=2.6in,width=3.2in]{../Figures/BadDrawing-InvSine}}}
    \centerline{The World's Most Evil Mathematician}
        \centerline{from {Math
            With Bad Drawings}}
          \centerline{by Ben Orlin}

When we were looking to restrict the range of the $\inverse\tan(x)$
there was a natural choice that was clearly visible in the
graph. Notice that the derivative of $\arctan(x)$ (and thus the
derivative of $\inverse\tan(x)$) is always positive while the
derivative of the $\arccot(x)$ (and thus the derivative of
$\inverse\cot(x)$) is always negative. This happens naturally for 
the tangent and cotangent, but not for the arcsine, the
arccosine or the other trigonometric functions.

%\pagebreak[4]
\begin{wrapfigure}[13]{r}{2in}
  \vskip-4mm{}
  \captionsetup{labelformat=empty}
  % \centerline{\includegraphics*[height=1.9in,width=1.3in]{../Figures/arcsin6}}
  \centerline{\includegraphics*[height=2.3in,width=2in]{../Figures/arcsin6}}
\label{fig:arcsin6}
\end{wrapfigure}
But this does suggest how we might proceed. As you can see in the
sketch at the right the slope of $\arcsin(x)$ is sometimes positive
and sometimes negative. Since we must restrict the range of
$\inverse\sin(x)$, we may as well make a convenient choice.  We'd like
to constrain the range of the inverse sine in such a way that its
derivative is always positive, and we'd like to constrain the range of
the inverse cosine in such a way that its derivative is always
negative.


From the sketch we see that if we restrict
the range of $\inverse\sin(x)$ to the interval $[-\pi/2, \pi/2]$ the
graph of $\inverse\sin(x)$ (in blue) will \emph{always}\aside{Except
  at the endpoints of course. What do you think the slope of
  $\inverse\sin(x)$ is
  at $x=\frac{\pi}{2}$?} have a
positive slope, and therefore a positive derivative.

With this restriction in place we \emph{define} the derivative of the
inverse sine as follows.

%\pagebreak[4]
\begin{mydefinition}[The Derivative of the Inverse Sine]
  \label{def:InvSin}
 The derivative of the inverse sine function is defined to be,\\
 \centerline{$\displaystyle\dfdx{(\inverse\sin(x))}{x}=\frac{1}{\sqrt{1-x^2}}$,}
 where $-\frac{\pi}{2}\lt \inverse\sin(x)\lt \frac{\pi}{2}$.
 % \vskip-4mm{}
\end{mydefinition}

\begin{embeddedproblem}
  What function is $\inverse\sin(x)$, as defined in
  Definition~\ref{def:InvSin} the inverse off?\\
  \hint{Naturally, $\inverse\sin(x)$ is the inverse of $\sin(x)$, but
    with which domain?}
\end{embeddedproblem}

Similarly, if we restrict the range of $\inverse\cos(x)$ to $0\lt x
\lt\pi$ its derivative will always be negative.

\begin{mydefinition}[The Derivative of the Inverse Cosine]
  \label{def:InvCos}
 The derivative of the inverse sine function is defined to be,\\
 \centerline{$\displaystyle\dfdx{(\inverse\cos(x))}{x}=\frac{-1}{\sqrt{1-x^2}}$,} where $0\lt x\lt \pi$.
%  \vskip-4mm{}
\end{mydefinition}
%\pagebreak[4]
% In Definition~\ref{def:InvSin} we had to choose a domain in order to
% ensure that $\inverse\sin(x)$ is a function. Many choices were
% possible so we made the convenient choice.  We will choose the domain
% of  $\inverse\cos(x)$ similarly.

% \begin{Digression}[Undefined Derivatives -- \textcolor{red}{Needs editing}]
%   % Notice that restricting the \emph{\underline{range}} of the inverse
%   % sine forces us to restrict the \emph{\underline{domain}} of the sine
%   % itself. Because they are mutually inverse the domain of one is the
%   % range of the other. So $\inverse\sin(x)$ is the inverse of $\sin(x)$
%   % with domain $[-\pi/2, \pi/2]$, but \alert{not} the inverse of
%   % $\sin(x)$ with domain $\RR$, because that function does not have an
%   % inverse.

%   % It is the nature of our topic which forces us to be so very careful
%   % in making our definitions, not any perverse desire for complexity
%   % for its own sake. Sometimes things are just complex and we have to
%   % deal with that.  But since we are having to pay such close attention
%   % to the details here we may as well point out one other small
%   % difficulty.
  
%   % If we don't think about it very hard it seems like the derivative of
%   % a function should be defined on the same domain as the function
%   % itself. Every point where a function is defined it should have a
%   % slope, right?

%   % But notice that this is not true of $\inverse\sin(x)$. The domain of
%   % $\inverse\sin(x)$, is $[-1,1]$ but the derivative of the inverse
%   % sine is clearly not defined at the endpoints. (Do you see why not?)
%   % So the domain of the inverse sine is slightly larger than the domain
%   % of its derivative.  Indeed, it is often the case that the domain of
%   % a function is (usually slightly) larger than the domain of its
%   % derivative.

  
%   We spend a good deal of time at the beginning of a Calculus course
%   talking about and drawing tangent lines simply because it is
%   helpful, and intuitively appealing, to do so. As a result of this
%   repetition it is easy to begin to identify the tangent line of the
%   graph of a function with the derivative of the function. But the
%   tangent line is not the derivative. They are two different entities.
  
  
%   They must be. Recall that in Definition~\ref{def:TangentLine} we
%   used the derivative to define the tangent line in
%   definition in the first place.

%   On the other hand when you look at the graph of $y=\inverse\sin(x)$
%   it seems intuitively clear that (vertical) tangent lines exist at
%   $x=-1$ and $x=1$.

%   But according to  Definition~\ref{def:TangentLine} they do not. The
%   problem is that since these lines are vertical they have no
%   slope. As a result neither
%   $\eval{\dfdx{\left(\inverse\sin(x)\right)}{x}}{x}{-1}$ nor
%   $\eval{\dfdx{\left(\inverse\sin(x)\right)}{x}}{x}{1}$ exists.

%   Since the derivative does not exist at $x=-1$, and $x=1$ it is
%   impossible for the tangent line, as we have defined it, to exist
%   because our definitions says that the slope of the tangent line at a
%   point is
%   equal to the derivative at that point.

%   This all probably seems like a lot of fussy, nit-picking on our
%   part. After all, we can \emph{see} in our graph that  \emph{tangent line} of $\inverse\sin(x)$ is
%   perfectly well defined at the endpoints. In both cases it is a
%   vertical line, so what's all the fuss about?
% has the same
%   slope a
%   So the
%   existence of the line tangent to the graph of $y=\inverse\sin(x)$ at
%   $x=\pm1$ doesn't necessarily mean that the derivative is defined
%   there. And in fact, in this case it is not.

%   Under what conditions a derivative is, and is not, defined is not a
%   matter we have taken up yet because we do not yet have all of the
%   tools we'll need. We will return to this question in
%   Section~\ref{subsec:underl-deriv}.
% \end{Digression}

% Inverting the cosine function has all of the same issues.  By flipping
% the graph of the cosine across the line $y=x$ we can see the graph of
% the multifunction, $\arccos(x)$. The cosine with domain $\RR$ is not
% invertible so inverting the cosine consists of selecting an
% appropriate branch of $\arccos(x)$ by constraining the the range of
% $\arccos(x)$ (equivalently, constraining the domain of $\cos(x)$).

% \begin{ProficiencyDrill}
%   \begin{wrapfigure}[]{r}{2in}
%     \vskip-5mm{}
% \captionsetup{labelformat=empty}
%   \centerline{\includegraphics*[height=1.4in,width=1.6in]{../Figures/arccos2}}
% \label{fig:arccos2}
% \end{wrapfigure}

% We used our freedom to constrain the range of $\inverse\sin(x)$ to
%   guarantee that $\dfdx{(\inverse\sin(x))}{x}$ is always
%   positive. Similarly we choose to constrain the range of $\inverse\cos(x)$ to
%   guarantee that $\dfdx{(\inverse\cos(x))}{x}$ is always negative.
%   Use the sketch at the right to show that if the
%   range is constrained to the interval $[0,\pi]$ then
%   $$
%   \dfdx{(\inverse\cos(x))}{x} = \frac{-1}{\sqrt{1-x^2}}.
%   $$
% \end{ProficiencyDrill}

%       %       \pagebreak[4]

\begin{embeddedproblem-enumerate}
  \begin{enumerate}[label={\bf{}(\alph*)}]
  \item Show that if $y=\arcsin(x)$ then $y$ satisfies the
    differential equation
    $$
    \dfdxn{y}{x}{2}+x\dfdx{y}{x}=0.
    $$
  \item The function   $y=\arccos(x)$ satisfies the same differential
    equation. Show this in two different ways.
    \begin{enumerate}[label={\bf{}(\roman*)}]{}
    \item By direct computation, just as you did part (a).
    \item By direct computation, after first observing that $\arccos(x)=\frac\pi2-\arcsin(x)$.
    \end{enumerate}
  \end{enumerate}{}
\end{embeddedproblem-enumerate}


      %       The derivative of inverse \emph{tangent} is always positive no matter which
      %       branch we choose. Similarly The derivative of
      %       inverse \alert{co}\emph{tangent} is always negative. There is no
      %       choice where the inverses of the tangent and cotangent are involved.

      %       But when we were choosing branches for $\inverse\sin(x)$ and
      %       $\inverse\cos(x)$ we found that we did have a choice. By choosing our
      %       branch carefully we can guarantee that either one has a positive or a
      %       negative derivative.  To be consistent with the inverse tangent we
      %       deliberately contrived to chose branches that made the derivative of
      %       the \emph{sine} positive and made the derivative of the
      %       \alert{co}\emph{sine} negative.

\subsection*{The Inverse Secant and Inverse Cosecant}
\begin{wrapfigure}[]{r}{2.5in}
  \vskip-6mm{}
  \captionsetup{labelformat=empty}
  % \centerline{\includegraphics*[height=1.5in,width=2in]{../Figures/InvSecTriangle}}
    \centerline{\includegraphics*[height=2in,width=2.5in]{../Figures/InvSecTriangle}}
  \label{fig:InvSecTriangle}
\end{wrapfigure}
To find the differential of $\arcsec(x)$ we refer to the sketch at the
right and we proceed exactly as before. If $y=\arcsec(x)$, then
$x=\sec(y)$ so
\begin{align*}
  \dx{(\sec(y))}      \amp{}=\dx{x}\\
  \sec(y)\tan(y)\dx{y}\amp{}= \dx{x}
                        \intertext{so that}
                        \dx{y}              \amp{}=\frac{{1}}{\sec(y)\tan(y)}\dx{x}.
\end{align*} 

\begin{ProficiencyDrill-1line}
  Show that: $ \dx\left(\arcsec(x)\right)=\frac{\pm1}{x\sqrt{x^2-1}}.  $
\end{ProficiencyDrill-1line}

%\pagebreak[4]
\begin{embeddedproblem-enumerate}{}
  \begin{enumerate}[label={\bf{}(\alph*)}]
  \item Use the fact that $y=\arccsc(x)$ if and only if $\csc(y)=x$
    and proceeding as we did above, show that
    $$
    \dx\left(\arccsc(x)\right)= \frac{\mp 1}{\sqrt{1-x^2}}\dx{x}.
    $$
  \item Derive the same result from the identity:  $\arccsc(x)=
    \frac{\pi}{2}-\arcsec(x)$.
  \end{enumerate}
\end{embeddedproblem-enumerate}


\begin{embeddedproblem-enumerate}
  \begin{enumerate}[label={\bf{}(\alph*)}]
  \item Show that if $y=\arcsec(x)$ then $y$ satisfies the
    differential equation
    $$
    x(x^2-1)\dfdxn{y}{x}{2}+(2x^2-1)\dfdx{y}{x}=0.
    $$
  \item The function   $y=\arccsc(x)$ satisfies the same differential
    equation. Show this in two different ways.
    \begin{enumerate}[label={\bf{}(\roman*)}]{}
    \item By direct computation, just as you did part (a).
    \item By direct computation, after first observing that $\arccsc(x)=\frac\pi2-\arcsec(x)$.
    \end{enumerate}
  \end{enumerate}{}
\end{embeddedproblem-enumerate}


Again we will want to define the inverse secant, $y=\inverse\sec(x)$
by restricting the range of the multifunction
$y=\arcsec(x)$ in such a way that the derivative of $\inverse\sec(x)$
will always be positive. Similarly for the inverse cosecant function.

But forcing the derivative of the inverse secant to always be positive
is a bit harder for a couple of reasons. The first is just that these
functions are used less and thus we are not as familiar with them.
The other is the nature of the derivative formula itself. Recall that
$$
\dfdx{(\arcsec(x))}{x}=\frac{\pm1}{x\sqrt{x^2-1}}.
$$
This time it is the presence of the $x$ in the denominator \emph{along
  with} the $\pm1$ in the numerator that is troublesome. By choosing
the range constraint judiciously we can control whether the plus or
the minus is chosen in the numerator but any reasonable constraint on
the \emph{range} of $\inverse\sec(x)$ will always include both
positive and negative numbers in its \emph{domain}. So the $x$ in the
denominator will make the derivative of $\inverse\sec(x)$ positive for
some values and negative for others. We don't want that. We want it to
be positive.

Usually when we want things in the world to be simple, Nature just
laughs at us. But not this time. This time the problem is simply that
we have too many choices.  Both the $\pm1$ in the numerator and
the $x$ in the denominator will influence the sign (positive or
negative) of the expression $\frac{\pm1}{x\sqrt{x^2-1}}$. If they are
both in play it is hard to control the sign. So to simplify things we
just declare, by fiat, that the $x$ in the denominator will be
positive.  Specifically, we'll replace the $x$ in the denominator
with $\abs{x}$. Now we can control the sign of the derivative of the
$\inverse\sec(x)$ by simply constraining its range.

\begin{embeddedproblem}
%   \begin{wrapfigure}[]{r}{2in}
%     \vskip-10mm{}
%     \captionsetup{labelformat=empty}
%     \centerline{\includegraphics*[height=1.6in,width=2in]{../Figures/arcsec5}}
%     \label{fig:arcsec5}
%   \end{wrapfigure}
   The sketch below \\
        \centerline{\includegraphics*[height=1.6in,width=2in]{../Figures/arcsec5}}
         shows that if we constrain the range of the
    $\inverse\sec(x)$ to be the interval $(0, \pi/2)\cup(\pi/2,\pi)$ then 
    $$
    \dfdx{(\inverse\sec(x))}{x}=\frac{1}{\abs{x}\sqrt{x^2-1}}.
    $$
  \begin{enumerate}[label={\bf{}(\alph*)}]
  \item  Use the identity $\inverse\csc(x)=\frac\pi2-\inverse\sec(x)$ to
    show that
    $
    \dfdx{(\inverse\csc(x))}{x}=\frac{-1}{\abs{x}\sqrt{x^2-1}}.
    $
%  \item What is the range of $\inverse\sec(x)$?
  \end{enumerate}
%  \vskip-3mm{}
\end{embeddedproblem}


% \begin{embeddedproblem-enumerate}
%   \begin{enumerate}[label={\bf{}(\alph*)}]
%   \item Sketch the graph of the multifunction $\arcsec(x)$ and use
%     it to identify the range constraint that guarantees that
%     $\dfdx{(\inverse\csc(x))}{x}$ is always negative.
%   \item What is the domain of $\inverse\sec(x)$?
%   \end{enumerate}
% \end{embeddedproblem-enumerate}

Until now there has been a pattern to differentiation that we have not
remarked on. But you may have noticed that when we differentiate a
function we usually get back another function of the same type.
That is,
\begin{enumerate}
\item The derivative of a polynomial is another polynomial,
\item The derivative of a quotient of polynomials is another quotient
  of  polynomials,
\item The derivative of a trigonometric function is another trigonometric
  function.
\end{enumerate}

The inverse trigonometric functions break this pattern. The derivative
of an inverse trigonometric function is not another inverse
trigonometric function. Nor is it a trigonometric function, a
polynomial, or a quotient of polynomials\aside{Well, actually that's
  not quite true the derivative of the inverse tangent \emph{is} a
  quotient of polynomials, and a particularly simple one at that. But
  it is not an inverse trigonometric function}.

Do you suppose there is any significance to this?





      %       when package did not graph
      %       $y=\arctan\left(\frac{1+x}{1-x}\right)$, but actually graphed
      %       $y=\inverse\tan\left(\frac{1+x}{1-x}\right)$.  So, what is the
      %       difference?  This brings up what we actually mean by a function.  In
      %       truth, $y=y(x)=\arctan(x)$ is not really a function.  To be a
      %       function, then for every $y$ value, there must be a single value that
      %       we can call $y(x)$. 

      %       Notice that this is \emph{not} the graph of a function.  This is
      %       because for any integer, $n$, $x=\tan(y)=\tan(x+n\pi)$.  
      %       Because of this restriction on the \emph{range} of $\tan(x)$ (and thus
      %       one the \emph{domain} of $\inverse\tan(x)$), when    $x\lt{}1$
      %       $\arctan\left(\frac{1+x}{1-x}\right)=\arctan(x)+\frac{\pi}{4}$.
      %       However, for $x\gt{}1$, $\arctan\left(\frac{1+x}{1-x}\right)=\arctan(x)-\frac{3\pi}{4}$.

      %       This can be clearly seen when we superimpose the graphs of
      %       $\arctan(x)+\frac{\pi}{4}$ (in orange) and
      %       $\arctan(x)-\frac{3\pi}{4}$ (in red)  on the two branches of
      %       $\arctan\left(\frac{1+x}{1-x}\right)$ (in black) as in the following sketch. \\
      %       \centerline{\includegraphics*[height=1.5in,width=4in]{../Figures/arctan4}}


      %       An easy way to remember what $\inverse{\tan}(x)$ means is to read
      %       the symbol $\inverse{\tan}(x)$ as ``the angle whose tangent is
      %       $x.$'' This emphasizes that $x$ is the tangent of some angle, and
      %       $\inverse{\tan}(x)$ is that angle.

      %       There similar issues with the ranges of the other five inverse
      %       trigonometric functions, $\inverse\sin,$ $\inverse\cos,$
      %       $\inverse\cot,$ $\inverse\sec,$ and $\inverse\csc$, but it is harder
      %       to see how to resolve them than it was for the inverse
      %       tangent. Because these problems consist mostly of defining
      %       consistent conventions for these functions we will just tell you
      %       what conventions are rather than discussing the conventions
      %       themselves at length as we have here.

      %       \begin{ProficiencyDrill}
      %       Because $\sec(x)=\frac{1}{\cos(x)}$ the
      %       restriction on the \emph{range} of the inverse secant function (and on
      %       the \emph{domain} of the secant) is the same as the restriction on the
      %       cosine and its inverse.

      %       Use this to show that
      %       $$
      %       \dx\left(\arcsec(x)\right)=\frac{1}{\abs{x}\sqrt{x^2-1}}.
      %       $$
      %       \hint{Don't let this problem frighten you. The absolute value
      %       follows directly from the \emph{domain} restriction. Use a triangle in
      %       the unit circle as before to get the formula.}
      %       \end{ProficiencyDrill}

%\pagebreak[4]

The inverse trigonometric functions will not have a significant role
in our story for a while. However they will be very useful in some
topics that will come up next semester. We've introduced them here
simply because it makes sense to develop all of the differentiation
rules at more or less the same time.

We now have  differentiation formulas for all of the inverse
trigonometric functions. These, along with the constraints on their
ranges are given in the table below.
\alert{Memorize them!} You will need them.

      %       \begin{table}[h]
      %       \caption{Differentials of Inverse Trigonometric functions}
      %       \label{tab:InvTrigDiff2}
\begin{center}
  \begin{tabular}{|cccc|} \hline Function\amp{}
                                           Domain\amp{}Range\amp{}Derivative\\\hline\hline $y=\inverse\tan(x)$ \amp{} Any
                                                                                                       real number \amp{} $-\frac{\pi}{2}\lt{}y\lt{}\frac{\pi}{2}$
                                                       \amp{}$\displaystyle\dfdx{y}{x}=\frac{1}{1+x^2}$\\
                                         \amp{}\amp{}\amp{}\\
    $y=\inverse\cot(x)$ \amp{} Any real number \amp{} $0\lt{}y\lt{}\pi$
                                                       \amp{}$\displaystyle\dfdx{y}{x}=\frac{-1}{1+x^2}$\\
                                         \amp{}\amp{}\amp{}\\
    $y=\inverse\sin(x)$ \amp{} $-1\le x\le 1$ \amp{}
                                           $-\frac{\pi}{2}\le
                                           y\le\frac{\pi}{2}$
                                                       \amp{}$\displaystyle\dfdx{y}{x}=\frac{1}{\sqrt{1-x^2}}$\\
                                         \amp{}\amp{}\amp{}\\
    $y=\inverse\cos(x)$ \amp{} $-1\le x\le 1$ \amp{} $0\le y\le\pi$
                                                       \amp{}$\displaystyle\dfdx{y}{x}=\frac{-1}{\sqrt{1-x^2}}$\\
                                         \amp{}\amp{}\amp{}\\
    \multirow{2}{*} {$y=\inverse\sec(x)$} \amp{}$x\le-1$
                                                 \amp{}$0\le y\lt{}\frac{\pi}{2}$ \amp{}\multirow{2}{*}
                                                                           {$\displaystyle\dfdx{y}{x}=\frac{1}{\abs{x}\sqrt{x^2-1}}$}\\
                                         \amp{} or $ x\ge 1$
                                                 \amp{}or $\frac{\pi}{2}\lt{}y\le\pi$\amp{}\\
                                         \amp{}\amp{}\amp{}\\
    \multirow{2}{*}{$y=\inverse\csc(x)$} \amp{}$x\le-1$
                                                 \amp{}$0\le y\lt{}\frac{\pi}{2}$
                                                       \amp{}\multirow{2}{*}{$\displaystyle\dfdx{y}{x}=\frac{-1}{\abs{x}\sqrt{x^2-1}}$}\\
                                         \amp{} or $ x\ge 1$
                                                 \amp{}or $-\frac{\pi}{2}\lt{}y\le0$\amp{}\\
                                         \amp{}\amp{}\amp{}\\\hline
  \end{tabular}
\end{center}
      %       \end{table}

%\pagebreak[4]
\begin{ProficiencyDrill}
  Compute each of the following:
  \begin{enumerate}[label={\bf{}(\alph*)}]
    \begin{multicols}{3}
    \item $\dx\left(\inverse\sin(x^{1/2})\right)$
    \item $\dx\left(\inverse\tan(x^{1/2})\right)$
    \item $\dx\left(\inverse\sec(x^{1/2})\right)$
    \item $\dx\left( \left( \inverse\sin(x) \right)^{-1}\right)$
    \item $\dx\left( \left( \inverse\tan(x) \right)^{-1}\right)$
    \item $\dx\left( \left( \inverse\sec(x) \right)^{-1}\right)$
    % \item
    %   $\dx \left[ (3x^2+2x-1) \inverse\sec(\inverse{x}+x^2) \right]$
    % \item $\dx \left(\inverse\csc(x^2) \right)$
    \end{multicols}
  \end{enumerate}
\end{ProficiencyDrill}

%\pagebreak[4]
\begin{embeddedproblem}[Find the Pattern]
  Show that for any constant $a\gt{}0$,
  \begin{enumerate}[label={\bf{}(\alph*)}]
    \begin{multicols}{2}
    \item
      $\dx{\left(\dfrac1a\inverse\tan\left(\dfrac{x}{a}\right)\right)}=\dfrac{1}{x^2+a^2}\dx{x}$
    \item
      $\dx{\left(\dfrac1a\inverse\sin\left(\dfrac{x}{a}\right)\right)}=\dfrac{1}{a\sqrt{a^2-x^2}}\dx{x}$
    \item
      $\dx{\left(\dfrac1a\inverse\sec\left(\dfrac{x}{a}\right)\right)}=\dfrac{1}{\abs{x}\sqrt{x^2-a^2}}\dx{x}$
    \item
      $\dx{\left(\dfrac1a\inverse\cot\left(\dfrac{a}{x}\right)\right)}=\dfrac{1}{x^2+a^2}\dx{x}$
    \item
      $\dx{\left(\dfrac1a\inverse\cos\left(\dfrac{a}{x}\right)\right)}=\dfrac{1}{\abs{x}\sqrt{x^2-a^2}}\dx{x}$
    \item
      $\dx{\left(\dfrac1a\inverse\csc\left(\dfrac{a}{x}\right)\right)}=\dfrac{1}{a\sqrt{a^2-x^2}}\dx{x}$
   \end{multicols}
 \end{enumerate}
 Identify and describe any patterns you see in the forms of these derivatives.
\end{embeddedproblem}

\begin{embeddedproblem}
  Show that
  $\displaystyle\dx{\left(\dfrac{\inverse\sin(x)+x\sqrt{1-x^2}}{2}\right)}=\sqrt{1-x^2}\dx{x}$.
\end{embeddedproblem}
%\vskip-2mm{}
\begin{embeddedproblem-enumerate}[Find the Pattern]
  \begin{enumerate}[label={\bf{}(\alph*)}]
    \item   Show that each of the following statements is true.
  \begin{enumerate}[label={\bf{}(\roman*)}]
    \begin{multicols}{2}
    \item
      $\dx\left(\inverse\sin\left(\frac{x}{\sqrt{x^2+1}}\right)\right)=\frac{1}{x^2+1}\dx{x}$
    \item
      $\dx\left(\inverse\sin\left(\frac{x}{\sqrt{x^2+4}}\right)\right)=\frac{2}{x^2+4}\dx{x}$
    \item
      $\dx\left(\inverse\sin\left(\frac{x}{\sqrt{x^2+9}}\right)\right)=\frac{3}{x^2+9}\dx{x}$
    \item
      $\dx\left(\inverse\sin\left(\frac{x}{\sqrt{x^2+5}}\right)\right)=\frac{\sqrt{5}}{x^2+5}\dx{x}$
    \end{multicols}
  \end{enumerate}
  \item Now compute
    $\dx\left(\inverse\sin\left(\frac{x}{\sqrt{x^2+a^2}}\right)\right),$
    where $a\gt0.$
  \item  Does this problem change if $a$ is negative? How?
  \end{enumerate}
%  \vskip-4mm{}
\end{embeddedproblem-enumerate}

%\vskip-2mm{}
%\pagebreak[4]
\begin{embeddedproblem}
  Show that each of the following statements is true.
  \begin{enumerate}[label={\bf{}(\alph*)}]
%    \begin{multicols}{2}
    \item
      $\dx\left(x\inverse\sin(x)+\sqrt{1-x^2}\right) =
      \inverse\sin(x)\dx{x}.$
    \item
      $\dx\left(x\inverse\cos(x)-\sqrt{1-x^2}\right) =
      \inverse\cos(x)\dx{x}.$
 %   \end{multicols}
    \end{enumerate}
%    \vskip-6mm{}
\end{embeddedproblem}
%\vskip-2mm{}
%\pagebreak[4]
\begin{embeddedproblem-enumerate}
  % Consider $y_1(x)=\inverse\cot\left(\frac{1-x}{1+x}\right),$
  % and $y_2(x)=\inverse\tan(x).$
  \begin{enumerate}[label={\bf{}(\alph*)}]
  \item Show that
    $\dx\left(\inverse\cot\left(\frac{1-x}{1+x}\right)\right)=
    \dx\left(\inverse\tan(x)\right)$.
  \item Part (a) says that the graphs of
    $y=\inverse\cot\left(\frac{1-x}{1+x}\right)$ and
    $y=\inverse\tan(x)$ have the same slope when $x\neq-1$  and that
    their graphs  should differ by a constant. Substitute $x=0$
    into both to see what this constant should be.
  \item Now plot the graphs of both. Do they differ by a constant?
    Explain why or why not.
    % \begin{multicols}{2}
  % \item Graph $y_1(x)$ and $y_2(x)$ on the same set of axes. You
  %   should see that when $x$ is positive they have \emph{nearly} the
  %   same graph.
  % \item Show that $\dx{y_1}$ and $\dx{y_2}$ are also \emph{nearly}
  %   the same, at least when $x\ge0.$
  % \item Modify $y_2$ so that it is identical to $y_1.$
    % \end{multicols}
  \end{enumerate}
\end{embeddedproblem-enumerate}

\section{Curvature}
\label{sec:curvature}
\TLogo{appendix:radi-curv-accord}

Intuitively speaking, the \term{curvature} of a plane curve is how
quickly a curve turns. It is an important property to understand when
designing, for example, a road or a roller-coaster where the sharpness
of a curve determines the safe speed limit for that curve. Also, the
stresses imposed on the physical structure of a moving object during
any high speed maneuver are directly related to the curvature of their
paths.

One of the first people to study the notion of curvature
was
\href{https://mathshistory.st-andrews.ac.uk/Biographies/Oresme/}{Nicole
  Oresme\aside{Pronouced ``Or-em,'' not ``Or-es-me.''}}
($1323-1382$). Oresme defined the curvature of a straight line to be
zero (no surprise), and he defined the curvature of a circle to be the
reciprocal of its radius. For example, a circle with radius $2$ is
$\frac12$ whereas the curvature of a circle with radius $1$ is
$1$. This makes intuitive sense. A circle with a
smaller radius is
``more curved'' that one with a larger radius.  The radius of the
earth is so large that the curvature of say, the equator, is so close
to zero as to be completely imperceptible, at least for someone
standing on the surface.

\begin{wrapfigure}[14]{r}{2.7in}
    \vskip-9mm{}
  \captionsetup{labelformat=empty}
  \centerline{\includegraphics*[height=2.5in,width=2.7in]{../Figures/Curvature6}}
  \label{fig:Curvature6}
\end{wrapfigure}
Oresme's definition of curvature works just fine for circles and
straight lines but we'd like to extend it to more general
curves. Specifically, given a curve, we'd like to associate a number
with each point on the curve which tells us the curvature at that
point. Ideally, this should also generalize Oresme's results for
straight lines and circles.
% should have a
% curvature of zero at every point and that the curvature of a circle
% should be be constantly equal to the reciprocal of the radius. The
% arctangent (multi)function is exactly the tool we need to begin.

The sketch at the right shows three circles in the first quadrant with
radii, $\frac13$, $1$, and $2$, respectively.  Notice that on the unit
circle, the point of tangency travels $\frac{\pi}{2}$ units along the
circumference while the the red tangent line rotates through an angle
of $\frac{\pi}{2}$ radians. On the circle of radius $2$, the point of
tangency travels $2 \times \frac{\pi}{2}$ units to turn the tangent
line the same $\frac{\pi}{2}$ radians, whereas on the circle of radius
$1/3$, the point of tangency travels $\frac13\times\frac{\pi}{2}$
units to turn the tangent line the same $\frac\pi2$ radians.  This
suggests that a possible definition for \term{curvature} might be the
number of radians the tangent line rotates through as it travels along
the curve divided by the distance traveled along the curve while
making the turn. In other words, the curvature for a circle of radius
$r$ would be
$$
\frac{\theta \text{ radians }}{r\theta \text{ units}}=\frac{1}{r}
\text{ radians per unit. }
$$
This seems to make intuitive sense too. A circle with a radius of $0.0000001$
meters certainly seems to be ``more curved'' than a circle with radius
$1,000,000,000$ meters.

By the same reasoning, the curvature of a straight line would be
zero. (Why?) 

This definition works fine for circles and straight lines which have
constant curvature, but what about a more general curve. What if the
curvature changes from point to point along the curve such as a long
and winding road.

Clearly we'll need a definition of curvature that also changes from point
to point on the curve. So it seems reasonable to consider infinitely
small changes in both the angle $\theta$ and in arclength $s$.  More
precisely, we make the following definition.


\begin{mydefinition}
  \label{def:curvature-1}
\begin{wrapfigure}[20]{l}{2in}
  \vskip-11mm{}
  \captionsetup{labelformat=empty}
  % \centerline{\includegraphics*[height=1.3in,width=1.3in]{../Figures/Curvature5}}
  \centerline{\includegraphics*[height=2in,width=2in]{../Figures/Curvature5}}
\label{fig:Curvature5}
\end{wrapfigure}
%\vskip-8mm{}
Suppose the  point $P=(x,y)$ lies 
on a curve $C$ as in the sketch at the left. Let $\theta$ 
represent
the angle formed in the figure at the
left by the line tangent to  curve $C$ at point $P$, with some
fixed line.

Then the curvature is defined to be
\begin{equation}
\kappa=\abs{\dfdx{\theta}{s}}.\label{eq:CurvatureDef}
\end{equation}
\end{mydefinition}

To keep things simple we will always designate $\theta$ to be the
angle the tangent line forms with a horizontal line.


It can be difficult at first to see how formula~\ref{eq:CurvatureDef}
measures curvature but think carefully for a moment about the meaning
of the expression $\dfdx{\theta}{s}$. Loosely put ``the rate of change
of $\theta$ (the angle of our curve with the horizontal) with respect
to the change of $s$ (the arclength of our curve)'' is the ratio of
how much a point moving on the curve has changed direction to how far
it has moved. The more our curve changes direction in a given distance
the more it is curved and so the more curvature it has.


Don't let the absolute value bars in this formula disturb you. They
are there simply because we don’t want to consider whether a curve
bends to the right or to the left (whether $\theta$ is increasing or
decreasing), but just how fast it bends.

Also note, for later reference, that
\begin{equation}
  \tan(\theta)=\dfdx{y}{x}.\label{eq:TanEqDeriv}
\end{equation}

% \begin{ProficiencyDrill}
%   Do you see that this definition captures the idea that curvature is
%   a measure of how quickly, or slowly, our direction changes as we
%   travel along the curve?  Explain.
% \end{ProficiencyDrill}

A good test of the value of a new definition is whether or not it
produces the same results that the earlier, more
intuitive musings that led you to it. It should be clear that with this definition,
the curvature of a straight line would still be zero. (Why?)

%\pagebreak[4]
\begin{embeddedproblem}
  \label{prob:CircleCurvature}
  Use the following sketch of a quarter-circle with radius $r$ to
  show that Definition~\ref{def:curvature-1} yields a curvature of $\kappa =\frac{1}{r}$.\\
  \centerline{\includegraphics*[height=2.1in,width=2.1in]{../Figures/Curvature3}}
\end{embeddedproblem}

In problem (\ref{prob:CircleCurvature}) it was helpful to know that the
line tangent to a circle is always perpendicular to its radius. But
for an arbitrary curve this is not necessarily true. Let's look at the
general situation.


First recall from equation~(\ref{eq:TanEqDeriv}) that
$ \tan(\theta) = \dfdx{y}{x}$.  Rewrite this as
$\theta=\arctan\left(\dfdx{y}{x}\right)$ and differentiate so that  
\begin{equation}
  \dx{\theta}=\left(\frac{1}{1+\left(\dfdx{y}{x}\right)^2}\right)\dx\hskip-2pt\left(\hskip-1mm\dfdx{y}{x}\right).\label{eq:DiffTheta}
\end{equation}

To continue we need to compute $\dx\left(\dfdx{y}{x}\right)$ which
seems straightforward enough since the Quotient Rule applies. This
results in
$$
\dx\left(\dfdx{y}{x}\right) =
\frac{\dx{x}\textcolor{red}{\dx{(\dx{y})}}-\dx{y}\textcolor{red}{\dx{(\dx{x})}}}{(\dx{x})^2}.
$$ 
Next simplify the notation a bit.  Let\aside{Note this very
  carefully, $\d{(\dx{y})}\neq (\dx{y})^2$,  because  the latter is
  $(\dx{y})\times(\dx{y})$ whereas the former is the \emph{second
    difference} of $y$.} $\textcolor{red}{\d{(\dx{y})}}=\dx^2{y}$ and
$\textcolor{red}{\d{(\dx{x})}}=\dx^2{x}$ so that,
$$
\dx\left(\dfdx{y}{x}\right) =
\frac{\dx{x}\d^2{y}-\dx{y}\dx^2{x}}{(\dx{x})^2}.
$$
% which seems OK at first glance. But there is a problem because
% $\dx^2x$ is equal to zero, but $\dx^2y$ is \emph{not necessarily}
% equal to zero and it is not at all clear why this should be true. We
% will take a moment to investigate this.

But $\dx(\dx{x}) =\dx^2x$ and $\dx(\dx{y}) =\dx^2y$ present a
problem for us.  We have observed regularly beginning  in
Chapter~\ref{chapt:differentials} that differentials are
questionable. If differentials are problematic, a
second differential -- the differential of a differential -- is even
moreso. The time has come to begin investigating this issue.

%\pagebreak[4]
\begin{Digression}[\foreign{Hic Sunt Dracones} (Here Be Dragons)]
  \label{digression:curvature-1}
\begin{wrapfigure}[13]{r}{1.5in}
  \vskip-4mm{}
\captionsetup{labelformat=empty}
\centerline{\href{https://en.wikipedia.org/wiki/Hunt-Lenox_Globe}{\includegraphics*[height=1.5in,width=1.5in]{../Figures/LennoxGlobe2}}}
\caption{From the  Rare Book Division of the The New York Public Library}
\label{fig:LennoxGlobe}
\end{wrapfigure}
To indicate that uncharted waters were the most dangerous medieval
mapmakers would fill in the unexplored regions of their maps with
illustrations of dragons and other mythological beasts. The Lennox
Globe (1510 CE, seen at the right) even has the inscription
``\foreign{Hic Sunt Dracones}'' (Here Be Dragons) along the eastern
Asian coast. 
  
The concept of the differential of a differential, like $\dx^2x$,- is a
mathematical dragon of sorts. % If we follow it too far we risk bringing
% all of  Calculus crashing to the ground in a fiery storm of logical
% contradictions.
  We will proceed carefully.

  The problem here is conceptual, not computational. That is, there is
  no computation we can do to see what the value of $\dx^2x$ should
  be. Rather, we need to think carefully about the meanings of our
  symbols.

  \subsubsection*{Leibniz' Conception}
  When we use Leibniz' differential notation we think of $\dx{y}$
  and $\dx{x}$ as tiny increments in the $x$ and $y$
  coordinates. Specifically, $\dx{x}=x_2-x_1$ where $x_2$ and $x_1$
  are \emph{infinitely} close together. As we've seen the concept of
  a differential is already a difficult thing to accept.

  But if we stick with Leibniz then it must be that
  $ \dx(\dx{x}) = \dx{x_2}-\dx{x_1}, $ right? We can write this down
  and we can read these symbols easily enough. But what they seem to
  \emph{mean} is that $\dx(\dx{x})$ is the infinitely small difference
  of two infinitely small differences.

   \vskip2mm{}
  \centerline{\bf{} Down that path there are dragons. We won't
    go there.} % \vskip2mm{}


  \subsubsection*{Newton's Conception}
  Instead we'll adopt Newton's viewpoint and hope that it gives us a
  meaningful interpretation of our symbols.
  
  For Newton, time was the only variable, and time always flows forward
  at a constant rate\aside{Of course, with Einstein's work on
    Relativity Theory we now know that this assumption is not true. No
    matter. Newton assumed it was true and he based his Method of
    Fluxions on that assumption.}. He thought of a variable, $x$ for
  example, as moving, or flowing (\foreign{fluent}) in time and
  $\dot{x}$ was the rate of flow (\foreign{fluxion}) of the
  \foreign{fluent} $x$. In the simplest example, if $x$ is the position
  of a point then $\dot{x}$ is its velocity.

  Newton only applied his dot notation to something that was changing
  in time; something that has a \emph{rate of change with respect to
    time.} Can we apply it to time itself?

  Sure we can. Time itself is ``changing in time'', isn't it? Time
  flows at a rate of $1$ second per second, or one day per day, or one
  century per century. The units don't really matter. Loosely speaking
  (and mixing the Newtonian and Leibnizian viewpoints a bit), an
  increment of time now has the same magnitude as an increment of time
  later, has the same magnitude as an increment of time in the
  past. That is, every increment of $t$ is the same size.  If we let
  $t$ represent elapsed time then this means that the infinitesimal
  increment $\dx{t}$ is \emph{constant}, so the fluxion of $\dx{t}$ is
  zero: $\dx(\dx{t})=\dx^2{t}=\dot{\dx{t}}=0$.  However, keep in mind
  that this is not true of a quantity that is
  flowing at a variable rate. For example, if $y=t^2$ then
  $\eval{\dx{y}}{t}{1}=2\dx{t}$ whereas $\eval{\dx{y}}{t}{2}=4\dx{t}$.

  This solves our problem. If $\dx{t}$ is \emph{constant} then by the
  Constant Rule $\dx^2{t}=\dx(\dx{t})=0$.

  Moreover we always have the option of adopting Newton's convention
  that the \emph{only} independent variable is time so that when we
  write $\dfdx{y}{x}$ we are thinking of  the independent
  variable $x$ as time (although we rarely say so out loud), and we are
  thinking of $y$ as functionally dependent on time (represented by
  $x$ in this case).  Therefore
  $$\dx^2x=\dx(\dx{x})=0 \text{ and } \dx^2y=\dx(\dx{y})\neq0.$$

  % Recall that the differential notation was Leibniz' invention and it
  % reflects a static viewpoint. But Newton's approach was dynamic. For
  % Newton the only \alert{independent} variable was time. Everything
  % else depends on how much time has elapse since we started our
  % clock. That is, everything else is a dependent variable (a
  % \term{function} of time).

  % Moreover Newton assumed that time flows forward into the future at a
  % \alert{constant} rate\aside{Of course, with Einstein's work on
  % Relativity Theory we now know that this assumption is not true.}
  % Since the differential of a
  % constant is zero we have, by the Constant Rule $\dx{(\dx{t})}=0$.


  As a practical matter we will rarely be interested in second order
  differentials like $\dx^2x$ or $\dx^2y$. As we indicated in
  Section~\ref{sec:hist-intr} differentials are really just an aid to
  computation; a convenient fiction. They are a helpful aid but the
  more important and useful concept is the derivative.  Differentials
  only serve as a step along the way to finding derivatives.

  But we will be very interested in second order derivatives:
  $
  \dfdx{}{x}\left(\dfdx{y}{x}\right) = \dfdxn{y}{x}{2}
  $
  because these often represent real, physical quantities like
  velocity or acceleration.
  

  But now there is another difficulty that can't be ignored.

  In order to make sense of the expression $\dx^2{x}$ we had to
  abandon Leibniz' conception altogether. It is nice that Newton's
  viewpoint gives us a way to make sense of the notation but what this
  tells us is that taking either viewpoint (Newton's or Leibniz') does
  not get us to the logical heart of the matter. A viewpoint is just a
  manner of thinking. We adopt a viewpoint in order to build
  intuition.

  When two viewpoints are at odds it means that we need something more
  general. In this case we need a theory that subsumes, and can
  replace, the views of both Leibniz and Newton.

  Eventually we will stand at the top of the Calculus mountain
  so we can see it all at once. From that vantage point we'll be able
  to see the path that Newton took, and the path that Leibniz took.
  This issue must be addressed but for now we will continue to work
  intuitively.  We will wait to begin our trek to the summit in
  chapter~\ref{chapt:whats-wrong-with}. 
\end{Digression}



Returning to our computation in progress we have
$$
dx\left(\dfdx{y}{x}\right) =\frac{\dx{x}\cdot
   \d^2{y}-\dx{y}\cdot\CancelToRed{=0}{\d^2{x}}}{(\dx{x})^2} = \left(\dfdxn{y}{x}{2}\right)\dx{x}% \frac{\dx{x}\cdot\d^2{y}}{(\dx{x})^2}
 $$

% \begin{align*}
%   \dx\left(\dfdx{y}{x}\right) \amp{}=
%                                 \frac{\dx{x}\cdot
%                                 \d^2{y}-\dx{y}\cdot\CancelToRed{=0}{\d^2{x}}}{(\dx{x})^2} \\
%   \dx\left(\dfdx{y}{x}\right) \amp{}= \left(\dfdxn{y}{x}{2}\right)\dx{x}% \frac{\dx{x}\cdot\d^2{y}}{(\dx{x})^2}
    %   \end{align*}
  
  Returning to equation~(\ref{eq:DiffTheta}) we have
  $$
  \dx{\theta}=\left(\frac{1}{1+\left(\dfdx{y}{x}\right)^2}\right)\dx\hskip-2pt\left(\dfdx{y}{x}\right)=\left(\frac{1}{1+\left(\dfdx{y}{x}\right)^2}\right)\left(\dfdxn{y}{x}{2}\right)\dx{x}
  $$
% \begin{align*}
%   \dx{\theta}\amp{}=\left(\frac{1}{1+\left(\dfdx{y}{x}\right)^2}\right)\dx\hskip-2pt\left(\dfdx{y}{x}\right)\\
%              \amp{}=\left(\frac{1}{1+\left(\dfdx{y}{x}\right)^2}\right)\left(\dfdxn{y}{x}{2}\right)\dx{x}% \frac{\dx{x}\cdot\d^2{y}}{(\dx{x})^2}\\
% \end{align*}

%\pagebreak[4]
\begin{embeddedproblem}
Show that
  \begin{equation}
    \kappa=\frac{\abs{\dfdxn{y}{x}{2}}}{\left[1+\left(\dfdx{y}{x}\right)^2\right]^{\frac32}}.\label{eq:CurvatureDef2}
  \end{equation}
  \hint{$\dx{s}=\sqrt{(\dx{x})^2+(\dx{y})^2}$.}
%  \vskip-6mm{}
\end{embeddedproblem}

Since the equation of a straight line is $y=mx+b$ it is clear that
formula~(\ref{eq:CurvatureDef2}) recovers Oresme's assertion that the
curvature of a straight line must be zero. (Why?)

\begin{myexample}
  Let's apply formula~(\ref{eq:CurvatureDef2}) to a circle to see if
  this agrees with Oresme's definition of the curvature of a circle.
  Consider the circle of radius $r$ given by the graph of
  $x^2+y^2=r^2$.
  Differentiate to obtain $  2x\dx{x}+2y\dx{y}=0$ so that \\
  \centerline{$\displaystyle \dfdx{y}{x}=-\frac{x}{y}.  $}
  Differentiate both sides again to see that\\
\centerline{$\displaystyle\dx\left(\hskip-2pt\dfdx{y}{x}\right)=\frac{-y\dx{x}+x\dx{y}}{y^2}$}
or, after dividing both sides by $\dx{x}$\\
\centerline{$\displaystyle  \dfdxn{y}{x}{2}=\frac{-y+x\dfdx{y}{x}}{y^2}$.}



% \begin{ProficiencyDrill-enumerate}
%     \begin{enumerate}[label={\bf{}(\alph*)}]
%     \item   Use the  values for $\dfdx{y}{x}$ and $\dfdxn{y}{x}{2}$
%       just computed to show that for a circle of radius $r$, the
%       curvature $\kappa=1/r$.
%     \item Does this agree with the idea that as the radius of a circle
%       gets larger, the curvature gets smaller?
% \end{enumerate}
% \end{ProficiencyDrill-enumerate}

\begin{ProficiencyDrill}
  Use the  values for $\dfdx{y}{x}$ and $\dfdxn{y}{x}{2}$ we
  just computed to show that for a circle of radius $r$, the
  curvature $\kappa=1/r$.
  Does this agree with the idea that as the radius of a circle
  gets larger, the curvature gets smaller?
\end{ProficiencyDrill}
\end{myexample}

As we said earlier, the curvature of a general curve need not be
constant, but we can still apply formula~(\ref{eq:CurvatureDef2}).

\begin{embeddedproblem-enumerate}
  \begin{enumerate}[label={\bf{}(\alph*)}]{}
    % \begin{multicols}{4}
  \item Show that the curvature of the parabola $y=x^2$ is given by
    $ \displaystyle\kappa=\frac{2}{(1+4x^2)^{\frac32}}$.
  \item Where does the greatest curvature occur? Does this agree with
    what you can see on the graph of $y=x^2$?
    % \end{multicols}
  \end{enumerate}
\end{embeddedproblem-enumerate}

\begin{embeddedproblem}
Consider the ellipse given by $\frac{x^2}{a^2}+\frac{y^2}{b^2}=1$,
  $a\gt{}b\gt{}0$.\\
%   \begin{wrapfigure}[5]{r}{2.3in}
%   \vskip-21mm{}
% \captionsetup{labelformat=empty}
  \centerline{\centerline{\includegraphics*[height=1.4in,width=2.3in]{../Figures/Curvature7}}}
% \label{fig:Curvature7}
% \end{wrapfigure}
  \begin{enumerate}[label={\bf{}(\alph*)}]
  \item Show that the curvature is given by
    $ \kappa=\frac{a^4b^4}{(a^4y^2+b^4x^2)^{\frac32}}$.\\
    \hint{Rewrite
      $\frac{x^2}{a^2}+\frac{y^2}{b^2}=1$ as
      $b^2x^2+a^2y^2=a^2b^2$.
    }
  \item Looking at this graph it appears that the curvature at $(a,0)$ should be
    greater than at $(0,b)$. Use the formula you just derived to
    verify this.
  \item A circle is the special case of an ellipse when $a=b$. Use the
    formula from part (a) to compute the curvature of a circle.
  \end{enumerate}
%  \vskip-6mm{}
\end{embeddedproblem}
%\vskip-3mm{}

\begin{embeddedproblem}[Descartes' Method of Normals and Curvature]
  Recall that in exercise~\#\ref{prob:DescartesNormals2} of
  section~\ref{sec:descartes-normals} we used Descartes's Method of
  Normals to find the
    \begin{wrapfigure}[7]{r}{2.4in}
%    \vskip-4mm{}
    \captionsetup{labelformat=empty}
    % \centerline{\includegraphics*[height=1.5in,width=2in]{../Figures/DescartesCircle5}}
    \centerline{\includegraphics*[height=1.8in,width=2.4in]{../Figures/DescartesCircle5}}
    \label{fig:DescartesCircle3}
  \end{wrapfigure}
line tangent to the graph of $y=\sqrt{2x}$ at
  the point $(2,2)$ by finding a circle with its center on the
  $x$-axis which touches the graph exactly once, as seen in the 
  sketch at the right.
  \begin{enumerate}[label={\bf{}(\alph*)}]
  % \item Take a guess, do you think the circle and the parabola in this
  %   sketch have the same curvature at $(2,2)$?
  \item Compute the curvatures of the parabola and the circle seen in
    the diagram at the right at their point of intersection to show
    that they are not the same.
  \item Find the equation of the circle which passes through the point
    $(2,2)$, and also has  the same slope
    and curvature as $y=\sqrt{2x}$ at the point $(2,2)$. This is known as the \term{osculating circle},
      or ``kissing circle''.
  \item Show that the center of the osculating circle is on the line
    which passes through $(2,2)$ and $(3,0)$.
  \end{enumerate}
%  \vskip-6mm{}
\end{embeddedproblem}

We observed earlier that if we  define curvature using
equation~(\ref{eq:CurvatureDef}) then the curvature of a straight line
will be zero as Oresme said it should be. The converse is also
true. Specifically, if $\abs{\dfdx{\theta}{s}}=\kappa=0$ then $\theta$
must be a constant and so $\dfdx{y}{x}=\tan(\theta)$ must also be a
constant. (For simplicity we'll ignore the case where
$\theta=\pm\frac{\pi}{2}$.)

\begin{ProficiencyDrill}
  Assume that  $-\frac{\pi}{2}\lt \theta \lt
  \frac{\pi}{2}$ and that $\theta$ is constant. Complete the line of
  reasoning we began in the previous paragraph to show that
  $$
  y=  \left(\tan(\theta)\right)x+b
  $$
  for some number $b$. Explain how you know that this is this the
  equation of a line.
\end{ProficiencyDrill}

\begin{wrapfigure}[]{r}{2in}
\captionsetup{labelformat=empty}
\centerline{\includegraphics*[height=1.5in,width=1.5in]{../Figures/Curvature5}}
\label{fig:Curvature5Redux}
\end{wrapfigure}We've also shown that the curvature of a circle is
also constant (specifically, it is the reciprocal of the circle's
radius). The converse of this is also true\aside{But only if we
  confine our attention to curves in the plane. It turns out that
  there are many curves in three-dimensional space, (that is, curves
  that are not constrained to a two-dimensional plane) with constant
  curvature which are neither circles nor lines.}. This is shown in
much the same manner as the straight line but, naturally, the
computations are more complex.

To see this recall the diagram that we used to define curvature, shown
at the right. If we  suppose that $\abs{\dfdx{\theta}{s}}=\kappa$ then
we see that $\dfdx{\theta}{s}=\pm\kappa$. From our diagram
we also see that
$$
\sin(\theta)=\dfdx{y}{s}\text{,  and }\cos(\theta)=\dfdx{x}{s}.
$$
\begin{embeddedproblem-enumerate}
         \begin{enumerate}[label={\bf{}(\alph*)}]{}
         \item Complete the line of reasoning we began in the
           previous paragraph to show that
           $\dx{y}=\pm\frac1\kappa\sin(\theta)\dx{\theta}$ and
           $\dx{x}=\pm\frac1\kappa\cos(\theta)\dx{\theta}$.
         \item Next show that
           \begin{align*}
             x\amp=\pm\frac1\kappa\sin(\theta) +x_0\\
             y\amp=\mp\frac1\kappa\cos(\theta) +y_0
           \end{align*}
           for some constants $x_0$ and $y_0$.
         \item And finally, show that the point $(x,y)$ lies on a
           circle of radius $\frac1\kappa$. What is the equation of the circle?
       \end{enumerate}

\end{embeddedproblem-enumerate}


%%% Local Variables: 
%%% outline-minor-mode: t
%%% eval:(flyspell-mode)
%%% TeX-master: "DiffCalc"
%%% End: 

\message{ !name(DiffCalc.tex) !offset(-1958) }
